\Article{5}
It will now be easy to assign to the multipliers \(x,y,z\) etc. such values that the sum
\[
LLmm+L'L'm'm'+L''L''m''m''+\text{etc.}
\]
attains its smallest value. Clearly, for this purpose it is not necessary to know the absolute magnitudes of the mean errors \(m,m',m''\) etc.; their relative proportions suffice. We shall therefore introduce in their place the weights of the observations \(p,p',p''\) etc., that is, numbers inversely proportional to the squares \(mm,m'm',m''m''\) etc., taking the weight of an arbitrarily chosen observation as unity. The quantities \(x,y,z\) etc. must therefore be determined so that the polynomial in the indeterminates
\[
\begin{gathered}
\frac{(ax+by+cz+\text{etc.}+l)^2}{p}+\frac{(a'x+b'y+c'z+\text{etc.}+l')^2}{p'}\\[4pt]
{}+\frac{(a''x+b''y+c''z+\text{etc.}+l'')^2}{p''}+\text{etc.}
\end{gathered}
\]
attains its smallest value, which we shall suppose to occur for \textit{determinate} values \(x^0,y^0,z^0\) etc.

Introducing the following notation
\SourceLeaf{067}
\[
\begin{aligned}
\frac{aa}{p}+\frac{a'a'}{p'}+\frac{a''a''}{p''}+\text{etc.}&=[aa]\\[4pt]
\frac{ab}{p}+\frac{a'b'}{p'}+\frac{a''b''}{p''}+\text{etc.}&=[ab]\\[4pt]
\frac{ac}{p}+\frac{a'c'}{p'}+\frac{a''c''}{p''}+\text{etc.}&=[ac]\\[4pt]
\frac{bb}{p}+\frac{b'b'}{p'}+\frac{b''b''}{p''}+\text{etc.}&=[bb]\\[4pt]
\frac{bc}{p}+\frac{b'c'}{p'}+\frac{b''c''}{p''}+\text{etc.}&=[bc]\\[4pt]
\frac{cc}{p}+\frac{c'c'}{p'}+\frac{c''c''}{p''}+\text{etc.}&=[cc]
\end{aligned}
\]
\begin{center}etc. and also\end{center}
\[
\begin{aligned}
\frac{al}{p}+\frac{a'l'}{p'}+\frac{a''l''}{p''}+\text{etc.}&=[al]\\[4pt]
\frac{bl}{p}+\frac{b'l'}{p'}+\frac{b''l''}{p''}+\text{etc.}&=[bl]\\[4pt]
\frac{cl}{p}+\frac{c'l'}{p'}+\frac{c''l''}{p''}+\text{etc.}&=[cl]
\end{aligned}
\]
\begin{center}etc.\end{center}
the condition for a minimum clearly requires that
\[
\left.\begin{aligned}
0&=[aa]x^0+[ab]y^0+[ac]z^0+\text{etc.}+[al]\\
0&=[ab]x^0+[bb]y^0+[bc]z^0+\text{etc.}+[bl]\\
0&=[ac]x^0+[bc]y^0+[cc]z^0+\text{etc.}+[cl]\\
&\text{etc.}
\end{aligned}\right\}\ (1)
\]
After the quantities \(x^0,y^0,z^0\) etc. have been derived from these equations by elimination, we shall set
\[
\left.\begin{aligned}
ax^0+by^0+cz^0+\text{etc.}+l&=L\\
a'x^0+b'y^0+c'z^0+\text{etc.}+l'&=L'\\
a''x^0+b''y^0+c''z^0+\text{etc.}+l''&=L''\\
\text{etc.}\quad&
\end{aligned}\right\}\ (2)
\]
With this done, the function of the quantities \(v,v',v''\) etc. that is best suited to determining our unknown quantity and subject to the least uncertainty will be the one whose partial derivatives, in the particular case under consideration, have the values \(L,L',L''\) etc. respectively; and the weight of this determination, which we shall denote by \(P\), will be
\SourceLeaf{068}
\[
=\frac{1}{\dfrac{LL}{p}+\dfrac{L'L'}{p'}+\dfrac{L''L''}{p''}+\text{etc.}}\quad(3)
\]
or \(\frac1P\) will be the value of the polynomial given above for the system of values of the quantities \(x,y,z\) etc. that satisfies equations (1).

\Article{6}
In the preceding article we showed how to identify the function \(U\) that serves to give the most suitable determination of our unknown quantity; let us now see what \textit{value} the unknown quantity attains in this way. Let this value be denoted by \(K\); it therefore results from substituting in \(U\) the observed values of the quantities \(v,v',v''\) etc.; let the same substitution give the function \(u\) the value \(k\); finally, let \(\kappa\) be the true value of the unknown quantity, which would therefore result from the true values of the quantities \(v,v',v''\) etc. if we could substitute these in either \(U\) or \(u\). Thus we shall have
\[
\begin{aligned}
k&=\kappa+le+l'e'+l''e''+\text{etc.}\\
K&=\kappa+Le+L'e'+L''l''+\text{etc.}
\end{aligned}
\]
and consequently
\[
K=k(L-l)e+(L'-l')e'+(L''-l'')e''+\text{etc.}
\]
Substituting in this equation for \(L-l,\ L'-l',\ L''-l''\) etc. the values from (2), and setting
\[
\left.\begin{aligned}
ae+a'e'+a''e''+\text{etc.}&=\mathfrak A\\
be+b'e'+b''e''+\text{etc.}&=\mathfrak B\\
ce+c'e'+c''e''+\text{etc.}&=\mathfrak C
\end{aligned}\right\}\ (4)
\]
etc., we shall have
\[
K=k+\mathfrak A x^0+\mathfrak B y^0+\mathfrak C z^0\ \text{etc.}\qquad(5)
\]
We cannot, indeed, calculate the values of the quantities \(\mathfrak A,\mathfrak B,\mathfrak C\) etc. by formulas (4), since the errors \(e,e',e''\) etc. remain unknown; but it is immediately clear that they are simply the values of the functions \(X,Y,Z\) etc. that result when the observed values are substituted for \(v,v',v''\) etc. In this way the system of equations (1), (3), (5) gives the complete solution of our problem, since the remark we made at the end of Art. 2. about basing the calculation of the quantities \(l,l',l''\) etc. on the observed values of the quan\SourceLeaf{069}%
tities \(v,v',v''\) etc. clearly applies with equal force to the calculation of the quantities \(a,a',a''\) etc. \(b,b',b''\) etc. etc.

\clearpage
